% % % % \input{appendix_denormalization_notice}
% Requires in preamble:
%   \usepackage[most]{tcolorbox}
%   \tcbuselibrary{listings,breakable,skins}
%   \usepackage{listings}
%   \usepackage{booktabs}   % for summary table only
%   \usepackage{enumitem}   % optional: compact lists (or remove [nosep] below)

\lstdefinestyle{sqlappendix}{%
  language=SQL,
  basicstyle=\ttfamily\footnotesize,
  breaklines=true,
  breakindent=0pt,
  breakautoindent=false,
  columns=flexible,
  keepspaces=true,
  showstringspaces=false,
  tabsize=2,
  aboveskip=0pt,
  belowskip=0pt,
  xleftmargin=0pt,
  framexleftmargin=0pt,
  keywordstyle=\bfseries\color{blue!60!black},
  commentstyle=\itshape\color{black!55},
  stringstyle=\color{teal!60!black},
  morekeywords={EXISTS, IN},%
}

\newtcblisting{sqlwrong}{%
  listing only,
  listing options={style=sqlappendix},
  enhanced,
  breakable,
  colback=red!4,
  colframe=red!40,
  boxrule=0.4pt,
  arc=2pt,
  left=5pt,
  right=5pt,
  top=3pt,
  bottom=3pt,
  before skip=4pt,
  after skip=5pt,
  title={\sffamily\bfseries Wrong},
  fonttitle=\small,
  coltitle=black,
  attach boxed title to top left={yshift=-1mm, xshift=2mm},
  boxed title style={empty, size=minimal},%
}

\newtcblisting{sqlcorrect}{%
  listing only,
  listing options={style=sqlappendix},
  enhanced,
  breakable,
  colback=green!4,
  colframe=green!45!black,
  boxrule=0.4pt,
  arc=2pt,
  left=5pt,
  right=5pt,
  top=3pt,
  bottom=3pt,
  before skip=4pt,
  after skip=8pt,
  title={\sffamily\bfseries Correct},
  fonttitle=\small,
  coltitle=black,
  attach boxed title to top left={yshift=-1mm, xshift=2mm},
  boxed title style={empty, size=minimal},%
}

\newtcblisting{sqlneutral}{%
  listing only,
  listing options={style=sqlappendix},
  enhanced,
  breakable,
  colback=gray!8,
  colframe=gray!45,
  boxrule=0.4pt,
  arc=2pt,
  left=5pt,
  right=5pt,
  top=3pt,
  bottom=3pt,
  before skip=4pt,
  after skip=8pt,%
}

\newcommand{\denormquestion}[1]{\par\smallskip\noindent\textit{Question:}~#1\par}
\newcommand{\denormrule}{\par\smallskip\noindent\textbf{RULE:}\space}

\subsection{Denormalization Notice}
\label{app:denormalization_notice}

\begin{tcolorbox}[
  colback=gray!5,
  colframe=gray!20,
  arc=4pt,
  boxrule=0.5pt,
  breakable,
  left=8pt,
  right=8pt,
  top=8pt,
  bottom=8pt,
  nobeforeafter,
]
{\bfseries Schema Denormalization Notice}\par\medskip

The database schema provided is in a denormalized form.
Tables have been constructed by joining multiple normalized entities together, which means a single logical record (e.g.\ one race result, one driver) may appear across multiple rows due to join redundancy.
You must account for this when writing SQL to avoid returning inflated or incorrect results.

The rules below tell you exactly how to handle this for each query type.

\tcbline
\paragraph{1.\ Retrieval Queries (SELECT without aggregation)}
Retrieval queries return rows of values---names, IDs, descriptions, dates.

\denormrule Always use \texttt{SELECT DISTINCT} when retrieving any column or combination of columns, unless the question explicitly asks for all occurrences including repetitions.

\noindent\textbf{Examples:}

\denormquestion{``List the names of all drivers who competed in race~5.''}
\begin{sqlwrong}
SELECT driver_name FROM results_flat WHERE race_id = 5
\end{sqlwrong}
\begin{sqlcorrect}
SELECT DISTINCT driver_name FROM results_flat WHERE race_id = 5
\end{sqlcorrect}

\denormquestion{``What circuits have hosted a race?''}
\begin{sqlwrong}
SELECT circuit_name FROM results_flat
\end{sqlwrong}
\begin{sqlcorrect}
SELECT DISTINCT circuit_name FROM results_flat
\end{sqlcorrect}

\noindent\textbf{When to \emph{not} use \texttt{DISTINCT} on retrieval:}
\begin{itemize}
  \item The question asks for all rows explicitly, e.g.\ ``list every lap time recorded''---repetition may be expected.
  \item The question includes \texttt{ORDER BY} with \texttt{LIMIT} and you are retrieving a specific ranked row---\texttt{DISTINCT} may change ranking behaviour.
\end{itemize}

\tcbline
\paragraph{2.\ COUNT Queries}
COUNT queries count how many things exist.
This is the most error-prone query type on denormalized schemas.
You must identify what entity is being counted and apply \texttt{DISTINCT} to that entity's primary identifier.

\denormrule Always use \texttt{COUNT(DISTINCT \textit{entity\_id})} rather than \texttt{COUNT(*)} or \texttt{COUNT(\textit{column})}, where \textit{entity\_id} is the primary key of the logical entity the question is asking about.

\noindent\textbf{Examples:}

\denormquestion{``How many drivers competed in race~5?''}
\begin{sqlwrong}
SELECT COUNT(*) FROM results_flat WHERE race_id = 5
\end{sqlwrong}
\begin{sqlwrong}
SELECT COUNT(driver_id) FROM results_flat WHERE race_id = 5
\end{sqlwrong}
\begin{sqlcorrect}
SELECT COUNT(DISTINCT driver_id) FROM results_flat WHERE race_id = 5
\end{sqlcorrect}

\denormquestion{``How many races has Hamilton participated in?''}
\begin{sqlwrong}
SELECT COUNT(*) FROM results_flat WHERE driver_name = 'Hamilton'
\end{sqlwrong}
\begin{sqlcorrect}
SELECT COUNT(DISTINCT race_id) FROM results_flat
  WHERE driver_name = 'Hamilton'
\end{sqlcorrect}

\denormquestion{``How many results were recorded with more than 5 points?''}
\begin{sqlcorrect}
SELECT COUNT(DISTINCT result_id) FROM results_flat WHERE points > 5
\end{sqlcorrect}

\noindent\textbf{Identifying which \textit{entity\_id} to use:}
\begin{itemize}
  \item If counting people/drivers/teams $\rightarrow$ \texttt{DISTINCT} on \texttt{driver\_id} or \texttt{team\_id}
  \item If counting events/races/games $\rightarrow$ \texttt{DISTINCT} on \texttt{race\_id} or \texttt{event\_id}
  \item If counting records/results $\rightarrow$ \texttt{DISTINCT} on \texttt{result\_id}
  \item If counting laps $\rightarrow$ lap rows are the atomic unit; \texttt{COUNT(*)} may be appropriate if laps themselves are what is counted
\end{itemize}

\tcbline
\paragraph{3.\ SUM and AVG Queries}
\texttt{SUM} and \texttt{AVG} are the most dangerous query types on denormalized schemas.
A simple \texttt{SUM(points)} will add up points once per row, but each logical result may appear many times (once per lap), inflating the total massively.
\texttt{DISTINCT} cannot be applied directly inside \texttt{SUM} or \texttt{AVG}.

\denormrule Always deduplicate using a subquery before applying \texttt{SUM} or \texttt{AVG}.
The subquery should \texttt{SELECT DISTINCT} on the primary key of the entity whose attribute you are aggregating, along with the attribute column.

\noindent\textbf{Template:}
\begin{sqlneutral}
SELECT SUM(col) / AVG(col)
FROM (
  SELECT DISTINCT entity_id, col
  FROM table_name
  WHERE condition
)
\end{sqlneutral}

\noindent\textbf{Examples:}

\denormquestion{``What is the total points scored by Hamilton?''}
\begin{sqlwrong}
SELECT SUM(points) FROM results_flat
  WHERE driver_name = 'Hamilton'
  -- Inflated: sums points once per lap row
\end{sqlwrong}
\begin{sqlcorrect}
SELECT SUM(points) FROM (
  SELECT DISTINCT result_id, points
  FROM results_flat
  WHERE driver_name = 'Hamilton'
)
\end{sqlcorrect}

\denormquestion{``What is the average points per race in the 2020 season?''}
\begin{sqlcorrect}
SELECT AVG(points) FROM (
  SELECT DISTINCT result_id, points
  FROM results_flat
  WHERE race_year = 2020
)
\end{sqlcorrect}

\tcbline
\paragraph{4.\ GROUP BY Queries}
\texttt{GROUP BY} queries aggregate within groups---e.g.\ total points per driver, number of races per year.
On denormalized schemas, \texttt{GROUP BY} alone does not deduplicate; it groups all rows including duplicates.

\denormrule When using \texttt{GROUP BY}, either:
\begin{enumerate}
  \item Apply \texttt{COUNT(DISTINCT \textit{entity\_id})} inside the aggregation, or
  \item Deduplicate with a subquery before grouping
\end{enumerate}

\noindent\textbf{Examples:}

\denormquestion{``How many races did each driver compete in?''}
\begin{sqlwrong}
SELECT driver_name, COUNT(*) FROM results_flat GROUP BY driver_name
\end{sqlwrong}
\begin{sqlcorrect}
SELECT driver_name, COUNT(DISTINCT race_id)
FROM results_flat
GROUP BY driver_name
\end{sqlcorrect}

\denormquestion{``What is the total points per driver in the 2020 season?''}
\begin{sqlwrong}
SELECT driver_name, SUM(points) FROM results_flat
  WHERE race_year = 2020
GROUP BY driver_name
\end{sqlwrong}
\begin{sqlcorrect}
SELECT driver_name, SUM(points) FROM (
  SELECT DISTINCT result_id, driver_name, points
  FROM results_flat
  WHERE race_year = 2020
) GROUP BY driver_name
\end{sqlcorrect}

\tcbline
\paragraph{5.\ MIN and MAX Queries}
\texttt{MIN} and \texttt{MAX} are safe on denormalized schemas because taking the minimum or maximum of a repeated value returns the same result regardless of how many times the value appears.

\denormrule No special handling needed for \texttt{MIN} and \texttt{MAX}. Use them normally.

\denormquestion{``What is the fastest lap time recorded by Hamilton?''}
\begin{sqlcorrect}
SELECT MIN(laptime) FROM laptimes_flat
  WHERE driver_name = 'Hamilton'
  -- Safe: MIN of repeated values is unaffected by duplicates
\end{sqlcorrect}

\tcbline
\paragraph{6.\ EXISTS and IN Subqueries}
\texttt{EXISTS} and \texttt{IN} check for the presence of a value, not its count.
Duplicates do not affect correctness here.

\denormrule No special handling needed for \texttt{EXISTS} or \texttt{IN} subqueries.

\tcbline
\paragraph{Summary Reference}
\renewcommand{\arraystretch}{1.15}
\begin{tabular}{@{}ll@{}}
  \toprule
  \textbf{Query type} & \textbf{Rule} \\
  \midrule
  \texttt{SELECT}     & Always \texttt{DISTINCT} unless repetition is intended \\
  \texttt{COUNT}      & Always \texttt{COUNT(DISTINCT entity\_id)} \\
  \texttt{SUM}/\texttt{AVG} & Deduplicate in subquery first, then aggregate \\
  \texttt{GROUP BY}   & Use \texttt{COUNT(DISTINCT ...)} or subquery deduplication \\
  \texttt{MIN}/\texttt{MAX} & Safe---no special handling needed \\
  \texttt{EXISTS}/\texttt{IN} & Safe---no special handling needed \\
  \bottomrule
\end{tabular}

\end{tcolorbox}

% Requires in preamble:
%   \usepackage[most]{tcolorbox}
%   \tcbuselibrary{listings,breakable,skins}
%   \usepackage{listings}
%   \usepackage{booktabs}   % for summary table only
%   \usepackage{enumitem}   % optional: compact lists (or remove [nosep] below)

\lstdefinestyle{sqlappendix}{%
  language=SQL,
  basicstyle=\ttfamily\footnotesize,
  breaklines=true,
  breakindent=0pt,
  breakautoindent=false,
  columns=flexible,
  keepspaces=true,
  showstringspaces=false,
  tabsize=2,
  aboveskip=0pt,
  belowskip=0pt,
  xleftmargin=0pt,
  framexleftmargin=0pt,
  keywordstyle=\bfseries\color{blue!60!black},
  commentstyle=\itshape\color{black!55},
  stringstyle=\color{teal!60!black},
  morekeywords={EXISTS, IN},%
}

\newtcblisting{sqlwrong}{%
  listing only,
  listing options={style=sqlappendix},
  enhanced,
  breakable,
  colback=red!4,
  colframe=red!40,
  boxrule=0.4pt,
  arc=2pt,
  left=5pt,
  right=5pt,
  top=3pt,
  bottom=3pt,
  before skip=4pt,
  after skip=5pt,
  title={\sffamily\bfseries Wrong},
  fonttitle=\small,
  coltitle=black,
  attach boxed title to top left={yshift=-1mm, xshift=2mm},
  boxed title style={empty, size=minimal},%
}

\newtcblisting{sqlcorrect}{%
  listing only,
  listing options={style=sqlappendix},
  enhanced,
  breakable,
  colback=green!4,
  colframe=green!45!black,
  boxrule=0.4pt,
  arc=2pt,
  left=5pt,
  right=5pt,
  top=3pt,
  bottom=3pt,
  before skip=4pt,
  after skip=8pt,
  title={\sffamily\bfseries Correct},
  fonttitle=\small,
  coltitle=black,
  attach boxed title to top left={yshift=-1mm, xshift=2mm},
  boxed title style={empty, size=minimal},%
}

\newtcblisting{sqlneutral}{%
  listing only,
  listing options={style=sqlappendix},
  enhanced,
  breakable,
  colback=gray!8,
  colframe=gray!45,
  boxrule=0.4pt,
  arc=2pt,
  left=5pt,
  right=5pt,
  top=3pt,
  bottom=3pt,
  before skip=4pt,
  after skip=8pt,%
}

\newcommand{\denormquestion}[1]{\par\smallskip\noindent\textit{Question:}~#1\par}
\newcommand{\denormrule}{\par\smallskip\noindent\textbf{RULE:}\space}

\subsection{Denormalization Notice}
\label{app:denormalization_notice}

\begin{tcolorbox}[
  colback=gray!5,
  colframe=gray!20,
  arc=4pt,
  boxrule=0.5pt,
  breakable,
  left=8pt,
  right=8pt,
  top=8pt,
  bottom=8pt,
  nobeforeafter,
]
{\bfseries Schema Denormalization Notice}\par\medskip

The database schema provided is in a denormalized form.
Tables have been constructed by joining multiple normalized entities together, which means a single logical record (e.g.\ one race result, one driver) may appear across multiple rows due to join redundancy.
You must account for this when writing SQL to avoid returning inflated or incorrect results.

The rules below tell you exactly how to handle this for each query type.

\tcbline
\paragraph{1.\ Retrieval Queries (SELECT without aggregation)}
Retrieval queries return rows of values---names, IDs, descriptions, dates.

\denormrule Always use \texttt{SELECT DISTINCT} when retrieving any column or combination of columns, unless the question explicitly asks for all occurrences including repetitions.

\noindent\textbf{Examples:}

\denormquestion{``List the names of all drivers who competed in race~5.''}
\begin{sqlwrong}
SELECT driver_name FROM results_flat WHERE race_id = 5
\end{sqlwrong}
\begin{sqlcorrect}
SELECT DISTINCT driver_name FROM results_flat WHERE race_id = 5
\end{sqlcorrect}

\denormquestion{``What circuits have hosted a race?''}
\begin{sqlwrong}
SELECT circuit_name FROM results_flat
\end{sqlwrong}
\begin{sqlcorrect}
SELECT DISTINCT circuit_name FROM results_flat
\end{sqlcorrect}

\noindent\textbf{When to \emph{not} use \texttt{DISTINCT} on retrieval:}
\begin{itemize}
  \item The question asks for all rows explicitly, e.g.\ ``list every lap time recorded''---repetition may be expected.
  \item The question includes \texttt{ORDER BY} with \texttt{LIMIT} and you are retrieving a specific ranked row---\texttt{DISTINCT} may change ranking behaviour.
\end{itemize}

\tcbline
\paragraph{2.\ COUNT Queries}
COUNT queries count how many things exist.
This is the most error-prone query type on denormalized schemas.
You must identify what entity is being counted and apply \texttt{DISTINCT} to that entity's primary identifier.

\denormrule Always use \texttt{COUNT(DISTINCT \textit{entity\_id})} rather than \texttt{COUNT(*)} or \texttt{COUNT(\textit{column})}, where \textit{entity\_id} is the primary key of the logical entity the question is asking about.

\noindent\textbf{Examples:}

\denormquestion{``How many drivers competed in race~5?''}
\begin{sqlwrong}
SELECT COUNT(*) FROM results_flat WHERE race_id = 5
\end{sqlwrong}
\begin{sqlwrong}
SELECT COUNT(driver_id) FROM results_flat WHERE race_id = 5
\end{sqlwrong}
\begin{sqlcorrect}
SELECT COUNT(DISTINCT driver_id) FROM results_flat WHERE race_id = 5
\end{sqlcorrect}

\denormquestion{``How many races has Hamilton participated in?''}
\begin{sqlwrong}
SELECT COUNT(*) FROM results_flat WHERE driver_name = 'Hamilton'
\end{sqlwrong}
\begin{sqlcorrect}
SELECT COUNT(DISTINCT race_id) FROM results_flat
  WHERE driver_name = 'Hamilton'
\end{sqlcorrect}

\denormquestion{``How many results were recorded with more than 5 points?''}
\begin{sqlcorrect}
SELECT COUNT(DISTINCT result_id) FROM results_flat WHERE points > 5
\end{sqlcorrect}

\noindent\textbf{Identifying which \textit{entity\_id} to use:}
\begin{itemize}
  \item If counting people/drivers/teams $\rightarrow$ \texttt{DISTINCT} on \texttt{driver\_id} or \texttt{team\_id}
  \item If counting events/races/games $\rightarrow$ \texttt{DISTINCT} on \texttt{race\_id} or \texttt{event\_id}
  \item If counting records/results $\rightarrow$ \texttt{DISTINCT} on \texttt{result\_id}
  \item If counting laps $\rightarrow$ lap rows are the atomic unit; \texttt{COUNT(*)} may be appropriate if laps themselves are what is counted
\end{itemize}

\tcbline
\paragraph{3.\ SUM and AVG Queries}
\texttt{SUM} and \texttt{AVG} are the most dangerous query types on denormalized schemas.
A simple \texttt{SUM(points)} will add up points once per row, but each logical result may appear many times (once per lap), inflating the total massively.
\texttt{DISTINCT} cannot be applied directly inside \texttt{SUM} or \texttt{AVG}.

\denormrule Always deduplicate using a subquery before applying \texttt{SUM} or \texttt{AVG}.
The subquery should \texttt{SELECT DISTINCT} on the primary key of the entity whose attribute you are aggregating, along with the attribute column.

\noindent\textbf{Template:}
\begin{sqlneutral}
SELECT SUM(col) / AVG(col)
FROM (
  SELECT DISTINCT entity_id, col
  FROM table_name
  WHERE condition
)
\end{sqlneutral}

\noindent\textbf{Examples:}

\denormquestion{``What is the total points scored by Hamilton?''}
\begin{sqlwrong}
SELECT SUM(points) FROM results_flat
  WHERE driver_name = 'Hamilton'
  -- Inflated: sums points once per lap row
\end{sqlwrong}
\begin{sqlcorrect}
SELECT SUM(points) FROM (
  SELECT DISTINCT result_id, points
  FROM results_flat
  WHERE driver_name = 'Hamilton'
)
\end{sqlcorrect}

\denormquestion{``What is the average points per race in the 2020 season?''}
\begin{sqlcorrect}
SELECT AVG(points) FROM (
  SELECT DISTINCT result_id, points
  FROM results_flat
  WHERE race_year = 2020
)
\end{sqlcorrect}

\tcbline
\paragraph{4.\ GROUP BY Queries}
\texttt{GROUP BY} queries aggregate within groups---e.g.\ total points per driver, number of races per year.
On denormalized schemas, \texttt{GROUP BY} alone does not deduplicate; it groups all rows including duplicates.

\denormrule When using \texttt{GROUP BY}, either:
\begin{enumerate}
  \item Apply \texttt{COUNT(DISTINCT \textit{entity\_id})} inside the aggregation, or
  \item Deduplicate with a subquery before grouping
\end{enumerate}

\noindent\textbf{Examples:}

\denormquestion{``How many races did each driver compete in?''}
\begin{sqlwrong}
SELECT driver_name, COUNT(*) FROM results_flat GROUP BY driver_name
\end{sqlwrong}
\begin{sqlcorrect}
SELECT driver_name, COUNT(DISTINCT race_id)
FROM results_flat
GROUP BY driver_name
\end{sqlcorrect}

\denormquestion{``What is the total points per driver in the 2020 season?''}
\begin{sqlwrong}
SELECT driver_name, SUM(points) FROM results_flat
  WHERE race_year = 2020
GROUP BY driver_name
\end{sqlwrong}
\begin{sqlcorrect}
SELECT driver_name, SUM(points) FROM (
  SELECT DISTINCT result_id, driver_name, points
  FROM results_flat
  WHERE race_year = 2020
) GROUP BY driver_name
\end{sqlcorrect}

\tcbline
\paragraph{5.\ MIN and MAX Queries}
\texttt{MIN} and \texttt{MAX} are safe on denormalized schemas because taking the minimum or maximum of a repeated value returns the same result regardless of how many times the value appears.

\denormrule No special handling needed for \texttt{MIN} and \texttt{MAX}. Use them normally.

\denormquestion{``What is the fastest lap time recorded by Hamilton?''}
\begin{sqlcorrect}
SELECT MIN(laptime) FROM laptimes_flat
  WHERE driver_name = 'Hamilton'
  -- Safe: MIN of repeated values is unaffected by duplicates
\end{sqlcorrect}

\tcbline
\paragraph{6.\ EXISTS and IN Subqueries}
\texttt{EXISTS} and \texttt{IN} check for the presence of a value, not its count.
Duplicates do not affect correctness here.

\denormrule No special handling needed for \texttt{EXISTS} or \texttt{IN} subqueries.

\tcbline
\paragraph{Summary Reference}
\renewcommand{\arraystretch}{1.15}
\begin{tabular}{@{}ll@{}}
  \toprule
  \textbf{Query type} & \textbf{Rule} \\
  \midrule
  \texttt{SELECT}     & Always \texttt{DISTINCT} unless repetition is intended \\
  \texttt{COUNT}      & Always \texttt{COUNT(DISTINCT entity\_id)} \\
  \texttt{SUM}/\texttt{AVG} & Deduplicate in subquery first, then aggregate \\
  \texttt{GROUP BY}   & Use \texttt{COUNT(DISTINCT ...)} or subquery deduplication \\
  \texttt{MIN}/\texttt{MAX} & Safe---no special handling needed \\
  \texttt{EXISTS}/\texttt{IN} & Safe---no special handling needed \\
  \bottomrule
\end{tabular}

\end{tcolorbox}

% Requires in preamble:
%   \usepackage[most]{tcolorbox}
%   \tcbuselibrary{listings,breakable,skins}
%   \usepackage{listings}
%   \usepackage{booktabs}   % for summary table only
%   \usepackage{enumitem}   % optional: compact lists (or remove [nosep] below)

\lstdefinestyle{sqlappendix}{%
  language=SQL,
  basicstyle=\ttfamily\footnotesize,
  breaklines=true,
  breakindent=0pt,
  breakautoindent=false,
  columns=flexible,
  keepspaces=true,
  showstringspaces=false,
  tabsize=2,
  aboveskip=0pt,
  belowskip=0pt,
  xleftmargin=0pt,
  framexleftmargin=0pt,
  keywordstyle=\bfseries\color{blue!60!black},
  commentstyle=\itshape\color{black!55},
  stringstyle=\color{teal!60!black},
  morekeywords={EXISTS, IN},%
}

\newtcblisting{sqlwrong}{%
  listing only,
  listing options={style=sqlappendix},
  enhanced,
  breakable,
  colback=red!4,
  colframe=red!40,
  boxrule=0.4pt,
  arc=2pt,
  left=5pt,
  right=5pt,
  top=3pt,
  bottom=3pt,
  before skip=4pt,
  after skip=5pt,
  title={\sffamily\bfseries Wrong},
  fonttitle=\small,
  coltitle=black,
  attach boxed title to top left={yshift=-1mm, xshift=2mm},
  boxed title style={empty, size=minimal},%
}

\newtcblisting{sqlcorrect}{%
  listing only,
  listing options={style=sqlappendix},
  enhanced,
  breakable,
  colback=green!4,
  colframe=green!45!black,
  boxrule=0.4pt,
  arc=2pt,
  left=5pt,
  right=5pt,
  top=3pt,
  bottom=3pt,
  before skip=4pt,
  after skip=8pt,
  title={\sffamily\bfseries Correct},
  fonttitle=\small,
  coltitle=black,
  attach boxed title to top left={yshift=-1mm, xshift=2mm},
  boxed title style={empty, size=minimal},%
}

\newtcblisting{sqlneutral}{%
  listing only,
  listing options={style=sqlappendix},
  enhanced,
  breakable,
  colback=gray!8,
  colframe=gray!45,
  boxrule=0.4pt,
  arc=2pt,
  left=5pt,
  right=5pt,
  top=3pt,
  bottom=3pt,
  before skip=4pt,
  after skip=8pt,%
}

\newcommand{\denormquestion}[1]{\par\smallskip\noindent\textit{Question:}~#1\par}
\newcommand{\denormrule}{\par\smallskip\noindent\textbf{RULE:}\space}

\subsection{Denormalization Notice}
\label{app:denormalization_notice}

\begin{tcolorbox}[
  colback=gray!5,
  colframe=gray!20,
  arc=4pt,
  boxrule=0.5pt,
  breakable,
  left=8pt,
  right=8pt,
  top=8pt,
  bottom=8pt,
  nobeforeafter,
]
{\bfseries Schema Denormalization Notice}\par\medskip

The database schema provided is in a denormalized form.
Tables have been constructed by joining multiple normalized entities together, which means a single logical record (e.g.\ one race result, one driver) may appear across multiple rows due to join redundancy.
You must account for this when writing SQL to avoid returning inflated or incorrect results.

The rules below tell you exactly how to handle this for each query type.

\tcbline
\paragraph{1.\ Retrieval Queries (SELECT without aggregation)}
Retrieval queries return rows of values---names, IDs, descriptions, dates.

\denormrule Always use \texttt{SELECT DISTINCT} when retrieving any column or combination of columns, unless the question explicitly asks for all occurrences including repetitions.

\noindent\textbf{Examples:}

\denormquestion{``List the names of all drivers who competed in race~5.''}
\begin{sqlwrong}
SELECT driver_name FROM results_flat WHERE race_id = 5
\end{sqlwrong}
\begin{sqlcorrect}
SELECT DISTINCT driver_name FROM results_flat WHERE race_id = 5
\end{sqlcorrect}

\denormquestion{``What circuits have hosted a race?''}
\begin{sqlwrong}
SELECT circuit_name FROM results_flat
\end{sqlwrong}
\begin{sqlcorrect}
SELECT DISTINCT circuit_name FROM results_flat
\end{sqlcorrect}

\noindent\textbf{When to \emph{not} use \texttt{DISTINCT} on retrieval:}
\begin{itemize}
  \item The question asks for all rows explicitly, e.g.\ ``list every lap time recorded''---repetition may be expected.
  \item The question includes \texttt{ORDER BY} with \texttt{LIMIT} and you are retrieving a specific ranked row---\texttt{DISTINCT} may change ranking behaviour.
\end{itemize}

\tcbline
\paragraph{2.\ COUNT Queries}
COUNT queries count how many things exist.
This is the most error-prone query type on denormalized schemas.
You must identify what entity is being counted and apply \texttt{DISTINCT} to that entity's primary identifier.

\denormrule Always use \texttt{COUNT(DISTINCT \textit{entity\_id})} rather than \texttt{COUNT(*)} or \texttt{COUNT(\textit{column})}, where \textit{entity\_id} is the primary key of the logical entity the question is asking about.

\noindent\textbf{Examples:}

\denormquestion{``How many drivers competed in race~5?''}
\begin{sqlwrong}
SELECT COUNT(*) FROM results_flat WHERE race_id = 5
\end{sqlwrong}
\begin{sqlwrong}
SELECT COUNT(driver_id) FROM results_flat WHERE race_id = 5
\end{sqlwrong}
\begin{sqlcorrect}
SELECT COUNT(DISTINCT driver_id) FROM results_flat WHERE race_id = 5
\end{sqlcorrect}

\denormquestion{``How many races has Hamilton participated in?''}
\begin{sqlwrong}
SELECT COUNT(*) FROM results_flat WHERE driver_name = 'Hamilton'
\end{sqlwrong}
\begin{sqlcorrect}
SELECT COUNT(DISTINCT race_id) FROM results_flat
  WHERE driver_name = 'Hamilton'
\end{sqlcorrect}

\denormquestion{``How many results were recorded with more than 5 points?''}
\begin{sqlcorrect}
SELECT COUNT(DISTINCT result_id) FROM results_flat WHERE points > 5
\end{sqlcorrect}

\noindent\textbf{Identifying which \textit{entity\_id} to use:}
\begin{itemize}
  \item If counting people/drivers/teams $\rightarrow$ \texttt{DISTINCT} on \texttt{driver\_id} or \texttt{team\_id}
  \item If counting events/races/games $\rightarrow$ \texttt{DISTINCT} on \texttt{race\_id} or \texttt{event\_id}
  \item If counting records/results $\rightarrow$ \texttt{DISTINCT} on \texttt{result\_id}
  \item If counting laps $\rightarrow$ lap rows are the atomic unit; \texttt{COUNT(*)} may be appropriate if laps themselves are what is counted
\end{itemize}

\tcbline
\paragraph{3.\ SUM and AVG Queries}
\texttt{SUM} and \texttt{AVG} are the most dangerous query types on denormalized schemas.
A simple \texttt{SUM(points)} will add up points once per row, but each logical result may appear many times (once per lap), inflating the total massively.
\texttt{DISTINCT} cannot be applied directly inside \texttt{SUM} or \texttt{AVG}.

\denormrule Always deduplicate using a subquery before applying \texttt{SUM} or \texttt{AVG}.
The subquery should \texttt{SELECT DISTINCT} on the primary key of the entity whose attribute you are aggregating, along with the attribute column.

\noindent\textbf{Template:}
\begin{sqlneutral}
SELECT SUM(col) / AVG(col)
FROM (
  SELECT DISTINCT entity_id, col
  FROM table_name
  WHERE condition
)
\end{sqlneutral}

\noindent\textbf{Examples:}

\denormquestion{``What is the total points scored by Hamilton?''}
\begin{sqlwrong}
SELECT SUM(points) FROM results_flat
  WHERE driver_name = 'Hamilton'
  -- Inflated: sums points once per lap row
\end{sqlwrong}
\begin{sqlcorrect}
SELECT SUM(points) FROM (
  SELECT DISTINCT result_id, points
  FROM results_flat
  WHERE driver_name = 'Hamilton'
)
\end{sqlcorrect}

\denormquestion{``What is the average points per race in the 2020 season?''}
\begin{sqlcorrect}
SELECT AVG(points) FROM (
  SELECT DISTINCT result_id, points
  FROM results_flat
  WHERE race_year = 2020
)
\end{sqlcorrect}

\tcbline
\paragraph{4.\ GROUP BY Queries}
\texttt{GROUP BY} queries aggregate within groups---e.g.\ total points per driver, number of races per year.
On denormalized schemas, \texttt{GROUP BY} alone does not deduplicate; it groups all rows including duplicates.

\denormrule When using \texttt{GROUP BY}, either:
\begin{enumerate}
  \item Apply \texttt{COUNT(DISTINCT \textit{entity\_id})} inside the aggregation, or
  \item Deduplicate with a subquery before grouping
\end{enumerate}

\noindent\textbf{Examples:}

\denormquestion{``How many races did each driver compete in?''}
\begin{sqlwrong}
SELECT driver_name, COUNT(*) FROM results_flat GROUP BY driver_name
\end{sqlwrong}
\begin{sqlcorrect}
SELECT driver_name, COUNT(DISTINCT race_id)
FROM results_flat
GROUP BY driver_name
\end{sqlcorrect}

\denormquestion{``What is the total points per driver in the 2020 season?''}
\begin{sqlwrong}
SELECT driver_name, SUM(points) FROM results_flat
  WHERE race_year = 2020
GROUP BY driver_name
\end{sqlwrong}
\begin{sqlcorrect}
SELECT driver_name, SUM(points) FROM (
  SELECT DISTINCT result_id, driver_name, points
  FROM results_flat
  WHERE race_year = 2020
) GROUP BY driver_name
\end{sqlcorrect}

\tcbline
\paragraph{5.\ MIN and MAX Queries}
\texttt{MIN} and \texttt{MAX} are safe on denormalized schemas because taking the minimum or maximum of a repeated value returns the same result regardless of how many times the value appears.

\denormrule No special handling needed for \texttt{MIN} and \texttt{MAX}. Use them normally.

\denormquestion{``What is the fastest lap time recorded by Hamilton?''}
\begin{sqlcorrect}
SELECT MIN(laptime) FROM laptimes_flat
  WHERE driver_name = 'Hamilton'
  -- Safe: MIN of repeated values is unaffected by duplicates
\end{sqlcorrect}

\tcbline
\paragraph{6.\ EXISTS and IN Subqueries}
\texttt{EXISTS} and \texttt{IN} check for the presence of a value, not its count.
Duplicates do not affect correctness here.

\denormrule No special handling needed for \texttt{EXISTS} or \texttt{IN} subqueries.

\tcbline
\paragraph{Summary Reference}
\renewcommand{\arraystretch}{1.15}
\begin{tabular}{@{}ll@{}}
  \toprule
  \textbf{Query type} & \textbf{Rule} \\
  \midrule
  \texttt{SELECT}     & Always \texttt{DISTINCT} unless repetition is intended \\
  \texttt{COUNT}      & Always \texttt{COUNT(DISTINCT entity\_id)} \\
  \texttt{SUM}/\texttt{AVG} & Deduplicate in subquery first, then aggregate \\
  \texttt{GROUP BY}   & Use \texttt{COUNT(DISTINCT ...)} or subquery deduplication \\
  \texttt{MIN}/\texttt{MAX} & Safe---no special handling needed \\
  \texttt{EXISTS}/\texttt{IN} & Safe---no special handling needed \\
  \bottomrule
\end{tabular}

\end{tcolorbox}

% Requires in preamble:
%   \usepackage[most]{tcolorbox}
%   \tcbuselibrary{listings,breakable,skins}
%   \usepackage{listings}
%   \usepackage{booktabs}   % for summary table only
%   \usepackage{enumitem}   % optional: compact lists (or remove [nosep] below)

\lstdefinestyle{sqlappendix}{%
  language=SQL,
  basicstyle=\ttfamily\footnotesize,
  breaklines=true,
  breakindent=0pt,
  breakautoindent=false,
  columns=flexible,
  keepspaces=true,
  showstringspaces=false,
  tabsize=2,
  aboveskip=0pt,
  belowskip=0pt,
  xleftmargin=0pt,
  framexleftmargin=0pt,
  keywordstyle=\bfseries\color{blue!60!black},
  commentstyle=\itshape\color{black!55},
  stringstyle=\color{teal!60!black},
  morekeywords={EXISTS, IN},%
}

\newtcblisting{sqlwrong}{%
  listing only,
  listing options={style=sqlappendix},
  enhanced,
  breakable,
  colback=red!4,
  colframe=red!40,
  boxrule=0.4pt,
  arc=2pt,
  left=5pt,
  right=5pt,
  top=3pt,
  bottom=3pt,
  before skip=4pt,
  after skip=5pt,
  title={\sffamily\bfseries Wrong},
  fonttitle=\small,
  coltitle=black,
  attach boxed title to top left={yshift=-1mm, xshift=2mm},
  boxed title style={empty, size=minimal},%
}

\newtcblisting{sqlcorrect}{%
  listing only,
  listing options={style=sqlappendix},
  enhanced,
  breakable,
  colback=green!4,
  colframe=green!45!black,
  boxrule=0.4pt,
  arc=2pt,
  left=5pt,
  right=5pt,
  top=3pt,
  bottom=3pt,
  before skip=4pt,
  after skip=8pt,
  title={\sffamily\bfseries Correct},
  fonttitle=\small,
  coltitle=black,
  attach boxed title to top left={yshift=-1mm, xshift=2mm},
  boxed title style={empty, size=minimal},%
}

\newtcblisting{sqlneutral}{%
  listing only,
  listing options={style=sqlappendix},
  enhanced,
  breakable,
  colback=gray!8,
  colframe=gray!45,
  boxrule=0.4pt,
  arc=2pt,
  left=5pt,
  right=5pt,
  top=3pt,
  bottom=3pt,
  before skip=4pt,
  after skip=8pt,%
}

\newcommand{\denormquestion}[1]{\par\smallskip\noindent\textit{Question:}~#1\par}
\newcommand{\denormrule}{\par\smallskip\noindent\textbf{RULE:}\space}

\subsection{Denormalization Notice}
\label{app:denormalization_notice}

\begin{tcolorbox}[
  colback=gray!5,
  colframe=gray!20,
  arc=4pt,
  boxrule=0.5pt,
  breakable,
  left=8pt,
  right=8pt,
  top=8pt,
  bottom=8pt,
  nobeforeafter,
]
{\bfseries Schema Denormalization Notice}\par\medskip

The database schema provided is in a denormalized form.
Tables have been constructed by joining multiple normalized entities together, which means a single logical record (e.g.\ one race result, one driver) may appear across multiple rows due to join redundancy.
You must account for this when writing SQL to avoid returning inflated or incorrect results.

The rules below tell you exactly how to handle this for each query type.

\tcbline
\paragraph{1.\ Retrieval Queries (SELECT without aggregation)}
Retrieval queries return rows of values---names, IDs, descriptions, dates.

\denormrule Always use \texttt{SELECT DISTINCT} when retrieving any column or combination of columns, unless the question explicitly asks for all occurrences including repetitions.

\noindent\textbf{Examples:}

\denormquestion{``List the names of all drivers who competed in race~5.''}
\begin{sqlwrong}
SELECT driver_name FROM results_flat WHERE race_id = 5
\end{sqlwrong}
\begin{sqlcorrect}
SELECT DISTINCT driver_name FROM results_flat WHERE race_id = 5
\end{sqlcorrect}

\denormquestion{``What circuits have hosted a race?''}
\begin{sqlwrong}
SELECT circuit_name FROM results_flat
\end{sqlwrong}
\begin{sqlcorrect}
SELECT DISTINCT circuit_name FROM results_flat
\end{sqlcorrect}

\noindent\textbf{When to \emph{not} use \texttt{DISTINCT} on retrieval:}
\begin{itemize}
  \item The question asks for all rows explicitly, e.g.\ ``list every lap time recorded''---repetition may be expected.
  \item The question includes \texttt{ORDER BY} with \texttt{LIMIT} and you are retrieving a specific ranked row---\texttt{DISTINCT} may change ranking behaviour.
\end{itemize}

\tcbline
\paragraph{2.\ COUNT Queries}
COUNT queries count how many things exist.
This is the most error-prone query type on denormalized schemas.
You must identify what entity is being counted and apply \texttt{DISTINCT} to that entity's primary identifier.

\denormrule Always use \texttt{COUNT(DISTINCT \textit{entity\_id})} rather than \texttt{COUNT(*)} or \texttt{COUNT(\textit{column})}, where \textit{entity\_id} is the primary key of the logical entity the question is asking about.

\noindent\textbf{Examples:}

\denormquestion{``How many drivers competed in race~5?''}
\begin{sqlwrong}
SELECT COUNT(*) FROM results_flat WHERE race_id = 5
\end{sqlwrong}
\begin{sqlwrong}
SELECT COUNT(driver_id) FROM results_flat WHERE race_id = 5
\end{sqlwrong}
\begin{sqlcorrect}
SELECT COUNT(DISTINCT driver_id) FROM results_flat WHERE race_id = 5
\end{sqlcorrect}

\denormquestion{``How many races has Hamilton participated in?''}
\begin{sqlwrong}
SELECT COUNT(*) FROM results_flat WHERE driver_name = 'Hamilton'
\end{sqlwrong}
\begin{sqlcorrect}
SELECT COUNT(DISTINCT race_id) FROM results_flat
  WHERE driver_name = 'Hamilton'
\end{sqlcorrect}

\denormquestion{``How many results were recorded with more than 5 points?''}
\begin{sqlcorrect}
SELECT COUNT(DISTINCT result_id) FROM results_flat WHERE points > 5
\end{sqlcorrect}

\noindent\textbf{Identifying which \textit{entity\_id} to use:}
\begin{itemize}
  \item If counting people/drivers/teams $\rightarrow$ \texttt{DISTINCT} on \texttt{driver\_id} or \texttt{team\_id}
  \item If counting events/races/games $\rightarrow$ \texttt{DISTINCT} on \texttt{race\_id} or \texttt{event\_id}
  \item If counting records/results $\rightarrow$ \texttt{DISTINCT} on \texttt{result\_id}
  \item If counting laps $\rightarrow$ lap rows are the atomic unit; \texttt{COUNT(*)} may be appropriate if laps themselves are what is counted
\end{itemize}

\tcbline
\paragraph{3.\ SUM and AVG Queries}
\texttt{SUM} and \texttt{AVG} are the most dangerous query types on denormalized schemas.
A simple \texttt{SUM(points)} will add up points once per row, but each logical result may appear many times (once per lap), inflating the total massively.
\texttt{DISTINCT} cannot be applied directly inside \texttt{SUM} or \texttt{AVG}.

\denormrule Always deduplicate using a subquery before applying \texttt{SUM} or \texttt{AVG}.
The subquery should \texttt{SELECT DISTINCT} on the primary key of the entity whose attribute you are aggregating, along with the attribute column.

\noindent\textbf{Template:}
\begin{sqlneutral}
SELECT SUM(col) / AVG(col)
FROM (
  SELECT DISTINCT entity_id, col
  FROM table_name
  WHERE condition
)
\end{sqlneutral}

\noindent\textbf{Examples:}

\denormquestion{``What is the total points scored by Hamilton?''}
\begin{sqlwrong}
SELECT SUM(points) FROM results_flat
  WHERE driver_name = 'Hamilton'
  -- Inflated: sums points once per lap row
\end{sqlwrong}
\begin{sqlcorrect}
SELECT SUM(points) FROM (
  SELECT DISTINCT result_id, points
  FROM results_flat
  WHERE driver_name = 'Hamilton'
)
\end{sqlcorrect}

\denormquestion{``What is the average points per race in the 2020 season?''}
\begin{sqlcorrect}
SELECT AVG(points) FROM (
  SELECT DISTINCT result_id, points
  FROM results_flat
  WHERE race_year = 2020
)
\end{sqlcorrect}

\tcbline
\paragraph{4.\ GROUP BY Queries}
\texttt{GROUP BY} queries aggregate within groups---e.g.\ total points per driver, number of races per year.
On denormalized schemas, \texttt{GROUP BY} alone does not deduplicate; it groups all rows including duplicates.

\denormrule When using \texttt{GROUP BY}, either:
\begin{enumerate}
  \item Apply \texttt{COUNT(DISTINCT \textit{entity\_id})} inside the aggregation, or
  \item Deduplicate with a subquery before grouping
\end{enumerate}

\noindent\textbf{Examples:}

\denormquestion{``How many races did each driver compete in?''}
\begin{sqlwrong}
SELECT driver_name, COUNT(*) FROM results_flat GROUP BY driver_name
\end{sqlwrong}
\begin{sqlcorrect}
SELECT driver_name, COUNT(DISTINCT race_id)
FROM results_flat
GROUP BY driver_name
\end{sqlcorrect}

\denormquestion{``What is the total points per driver in the 2020 season?''}
\begin{sqlwrong}
SELECT driver_name, SUM(points) FROM results_flat
  WHERE race_year = 2020
GROUP BY driver_name
\end{sqlwrong}
\begin{sqlcorrect}
SELECT driver_name, SUM(points) FROM (
  SELECT DISTINCT result_id, driver_name, points
  FROM results_flat
  WHERE race_year = 2020
) GROUP BY driver_name
\end{sqlcorrect}

\tcbline
\paragraph{5.\ MIN and MAX Queries}
\texttt{MIN} and \texttt{MAX} are safe on denormalized schemas because taking the minimum or maximum of a repeated value returns the same result regardless of how many times the value appears.

\denormrule No special handling needed for \texttt{MIN} and \texttt{MAX}. Use them normally.

\denormquestion{``What is the fastest lap time recorded by Hamilton?''}
\begin{sqlcorrect}
SELECT MIN(laptime) FROM laptimes_flat
  WHERE driver_name = 'Hamilton'
  -- Safe: MIN of repeated values is unaffected by duplicates
\end{sqlcorrect}

\tcbline
\paragraph{6.\ EXISTS and IN Subqueries}
\texttt{EXISTS} and \texttt{IN} check for the presence of a value, not its count.
Duplicates do not affect correctness here.

\denormrule No special handling needed for \texttt{EXISTS} or \texttt{IN} subqueries.

\tcbline
\paragraph{Summary Reference}
\renewcommand{\arraystretch}{1.15}
\begin{tabular}{@{}ll@{}}
  \toprule
  \textbf{Query type} & \textbf{Rule} \\
  \midrule
  \texttt{SELECT}     & Always \texttt{DISTINCT} unless repetition is intended \\
  \texttt{COUNT}      & Always \texttt{COUNT(DISTINCT entity\_id)} \\
  \texttt{SUM}/\texttt{AVG} & Deduplicate in subquery first, then aggregate \\
  \texttt{GROUP BY}   & Use \texttt{COUNT(DISTINCT ...)} or subquery deduplication \\
  \texttt{MIN}/\texttt{MAX} & Safe---no special handling needed \\
  \texttt{EXISTS}/\texttt{IN} & Safe---no special handling needed \\
  \bottomrule
\end{tabular}

\end{tcolorbox}
